\documentclass[11pt]{article}

\usepackage[T1]{fontenc}
\usepackage[utf8]{inputenc}
% \usepackage[spanish,es-nodecimaldot]{babel}
\usepackage{babel}
\usepackage{lmodern}
\usepackage[margin=1in]{geometry}
\usepackage{amsmath,amssymb}

\usepackage{xcolor}

\setlength{\parindent}{0pt}
\setlength{\parskip}{0.45em}
\pagestyle{plain}
\allowdisplaybreaks

\newcommand{\MRS}{\mathrm{MRS}}
\newcommand{\MPN}{\mathrm{MPN}}

\begin{document}

\begin{center}
  {\large\textbf{Centro de Investigación y Docencia Económicas}}\\[0.25em]
  {\large\textbf{Macroeconomía II}}\\[0.2em]
  {\large\textbf{Problem Set 2: Equilibrio General}}\\[0.2em]
  {\large\textbf{Solver}}\\[0.35em]
  LECO -- Otoño 2026
\end{center}

\vspace{0.6em}
\hrule
\vspace{0.8em}

\noindent\textbf{Ejercicio 1. Equilibrio general en una economía cerrada de un periodo}

La interacción entre hogares y empresas determina el equilibrio macroeconómico, el cuál se compone de una colección de variables endógenas que cumplen con las condiciones que recolectamos en la definición de un equilibrio competitivo.

\subsection*{(a) Definición del equilibrio competitivo}

Un \textbf{equilibrio competitivo} es una asignación
$
  \left(C^*,\ell^*, N^{d*}\right)
$, un sistema impositivo $T^*$ y un vector de precios $(1,w^*)$ tales que, dadas las variables exógenas $(z,G)$:\footnote{El precio del bien de consumo se normaliza a uno por ser el numerario. Por lo tanto, el vector de precios de equilibrio es $(1,w^*)$ y el salario real $w^*$ es el único precio relativo endógeno de esta economía.}
\begin{enumerate}
  \item $(C^*,\ell^*)$ resuelve el problema del hogar:
  \[
    \max_{C,\ell}
    \ U(C,\ell) = \log C+\psi\log\ell
    \quad 
    \text{s.a.}
    \quad 
    C\leq w^*N^s+\pi^*-T^*, 
    \quad N^s = 1 - \ell
  \]
  tomando $w^*$ como dado y $(\pi^*, T^*)$;

  \item $N^{d*}$ resuelve el problema de la empresa:
  \[
    \max_{N^d}
    \ \pi = Y-w^*N^d, 
        \quad 
    \text{s.a.}
    \quad 
    Y = z\sqrt{N^d}
  \]
  tomando $w^*$ como dado;
  \item El gobierno cumple con su restricción presupuestaria $G = T^*$;
  \item El mercado laboral y el mercado de bienes se vacían:
  \[
    N^{s*} = N^{d*},
  \]
  \[
    Y^*=C^*+G,
  \]
  donde $N^{s*}=1-\ell^*$ y $Y^* = z\sqrt{N^{d*}}$.
\end{enumerate}



\subsection*{(b) Reglas de decisión del hogar y de la empresa}

En cursos anteriores estudiaron la macroeconomía imponiendo directamente relaciones entre variables agregadas. Por ejemplo, se postulaba que el consumo privado depende del ingreso disponible, \(C=C(Y^D)\), suponiendo además que \(\partial C/\partial Y^D>0\). Una especificación habitual de esta relación es una función lineal del tipo $C=c_0+c_1Y^D.$\footnote{Los parámetros \(c_0\) y \(c_1\) son coeficientes de una relación agregada en \textit{forma reducida}. Las \textit{ecuaciones en forma reducida} describen cómo se relacionan las variables observadas, pero no identifican directamente los parámetros y mecanismos económicos que la generan (e.g. parámetros asociados a las preferencias, tecnología u otros mecanismos económicos). Las \textit{ecuaciones estructurales}, en cambio, especifican esos mecanismos de manera explícita. Un sistema compuesto por ecuaciones estructurales constituye un \textit{modelo estructural}. Al resolver un modelo estructural, los coeficientes de su forma reducida resultan de la interacción entre sus parámetros estructurales. Una característica importante de las ecuaciones en forma reducida es que sus parámetros pueden cambiar cuando cambia la política económica, mientras que los parámetros estructurales son invariantes respecto a la política económica (\textit{policy-invariant}). El modelo de este ejercicio, y en general todos los modelos del curso, son modelos estructurales que en primera instancia son invariantes ante cambios en la política. Volveremos a esta discusión al estudiar la crítica de Lucas en el tema de ciclos económicos reales.}

En la macroeconomía microfundamentada buscamos explicar el consumo agregado a partir de las decisiones óptimas de los hogares, las cuáles dependen de sus preferencias y restricciones, en lugar de postular directamente una función agregada de consumo. En un modelo con varios hogares, el consumo agregado es la suma de esas elecciones óptimas individuales: \(C=\sum_i C_i\). En un modelo de agente representativo, el consumo agregado es simplemente la elección óptima del hogar representativo.



En este inciso buscamos expresar las decisiones óptimas de los agentes como funciones del salario real, de parámetros y de las variables exógenas del modelo. El salario real es el único precio relativo de esta economía y, por lo tanto, la única variable de ajuste. Las variables exógenas, denominadas \textit{shifters} en la literatura, son variables determinadas fuera del modelo que generan desplazamientos de las curvas. Por otro lado, los parámetros asociados a las funciones determinan características como curvatura, puntos críticos, etc.

Las funciones que buscamos, normalmente conocidas como \textit{reglas de decisión}, son entonces:
\[
\ell(w; z, G), \qquad C^d(w; z, G), \qquad  N^d(w; z, G),
\]
Para simplificar notación, es común no escribir explícitamente la dependencia en variables exógenas y dejarlas únicamente en función de la variable independiente de interés, que es el precio relativo $w$. Es decir,
$$
\ell(w), \qquad C^d(w), \qquad  N^d(w).
$$
A partir de estas reglas de decisión se construyen las funciones microfundamentadas de oferta y demanda de los mercados laboral y de bienes y servicios:
\[
N^s(w) = 1 - \ell(w),\qquad N^d(w),\qquad
Y^s(w) = z \sqrt{N^d(w)},\qquad Y^d(w)=C^d(w)+G.
\]

En el siguiente inciso utilizaremos estas funciones e impondremos simultáneamente las condiciones de vaciado de mercado para determinar el salario real de equilibrio \(w^*\) y, por consecuencia, la asignación de equilibrio.


\medskip

\subsubsection*{El problema del hogar}

Como la utilidad del hogar es estrictamente creciente en el consumo, la restricción presupuestaria se satisface con igualdad:
\[
C=w(1-\ell)+\pi-T.
\]

El lagrangiano asociado al problema del hogar es entonces
\[
\mathcal L(C,\ell,\lambda)
=
\log C+\psi\log\ell
+\lambda\left[w(1-\ell)+\pi-T-C\right].
\]

Las C.P.O. de una solución interior son:
\begin{align*}
\frac{\partial\mathcal L}{\partial C}
&=\frac{1}{C}-\lambda=0,\\
\frac{\partial\mathcal L}{\partial\ell}
&=\frac{\psi}{\ell}-\lambda w=0,\\
\frac{\partial\mathcal L}{\partial\lambda}
&=w(1-\ell)+\pi-T-C=0.
\end{align*}
Las C.P.O. forman un sistema de 3 ecuaciones con 3 incógnitas: $(C, \ell, \lambda)$.

De la primera condición se obtiene
\[
\lambda=\frac{1}{C}.
\]
Al sustituir esta expresión en la segunda condición,
\[
\frac{\psi}{\ell}=\frac{w}{C}.
\]
Por lo tanto,
\begin{equation}
  C=\frac{w\ell}{\psi}.\label{eq:c_1}
\end{equation}


Sustituyendo esta expresión en la tercera condición,
\[
\frac{w\ell}{\psi}=w(1-\ell)+\pi-T.
\]
Multiplicando ambos lados por \(\psi\),
\[
w\ell=\psi w(1-\ell)+\psi\pi-\psi T,
\]
de modo que
\[
w\ell=\psi w-\psi w\ell+\psi\pi-\psi T.
\]
Reuniendo los términos que contienen \(\ell\),
\[
(1+\psi)w\ell=\psi(w+\pi-T).
\]
Por lo tanto,
\begin{equation}
\ell(w,\pi,T)
=
\frac{\psi(w+\pi-T)}{(1+\psi)w}. \label{eq:ell_p}
\end{equation}

Sustituyendo \eqref{eq:ell_p} en \eqref{eq:c_1}:
\begin{equation}
C^d(w,\pi,T) = 
\frac{w+\pi-T}{1+\psi}. \label{eq:C_p}
\end{equation}

Sustituyendo \eqref{eq:C_p} en la oferta laboral,
\begin{align*}
N^s(w,\pi,T)
&=1-\ell(w,\pi,T)\\
&=1-\frac{\psi(w+\pi-T)}{(1+\psi)w}\\
&=\frac{(1+\psi)w-\psi(w+\pi-T)}
        {(1+\psi)w}\\
&=\frac{w-\psi\pi+\psi T}{(1+\psi)w}.
\end{align*}

Las ecuaciones \eqref{eq:ell_p}-\eqref{eq:C_p} son la solución del problema de elección de consumo-ocio del hogar en \textbf{equilibrio parcial}.\footnote{Es decir, son las funciones de demanda de ocio, demanda de consumo y oferta laboral en equilibrio parcial.} Estas describen las decisiones óptimas del hogar en función de \(w\), tomando \(\pi\) y \(T\) como dados exógenamente. Sin embargo, en equilibrio general, \(\pi\) y \(T\) son variables endógenas. En particular, el impuesto de suma fija se determina mediante la restricción presupuestaria del gobierno: $G = T.$ Por su parte, los beneficios se determinan endógenamente a partir del problema de maximización de la empresa. Por lo tanto, para encontrar las funciones de oferta laboral y demanda de bienes en equilibrio general, es necesario expresar a estas en función únicamente de $w$, de parámetros y de las variables exógenas \((z,G)\). Como $G=T$, y $G$ es una variable exógena, entonces basta con sustituir $T$ por $G$ en cada función y resolver el problema de la empresa para obtener \(\pi=\pi(w)\).

\subsubsection*{El problema de la empresa}

La empresa representativa elige su demanda de trabajo para maximizar
beneficios:
\[
\max_{N^d\geq 0}
\left\{
z\sqrt{N^d}-wN^d
\right\}.
\]

La condición de primer orden es
\[
\frac{z}{2\sqrt{N^d}}-w=0.
\]

Despejando \(N^d\),
\[
\frac{z}{2\sqrt{N^d}}=w,
\]
\[
z=2w\sqrt{N^d},
\]
\[
\sqrt{N^d}=\frac{z}{2w}.
\]
Por lo tanto, la \textcolor{blue}{\textbf{función de demanda laboral}} es
\begin{equation}
  \boxed{
N^d(w)=\frac{z^2}{4w^2}\label{eq:N_d}
}.
\end{equation}

La demanda laboral tiene pendiente negativa y es convexa respecto al salario real:
$$
N^d_w = -\frac{z^2}{2w^3}<0, \quad N^d_{ww} = \frac{3z^2}{2w^4}>0, \quad \forall w>0.
$$

La función de oferta de bienes se obtiene sustituyendo la demanda de trabajo en la función de producción:
\begin{align*}
Y^s(w)
&=z\sqrt{N^d(w)}\\
&=z\sqrt{\frac{z^2}{4w^2}}\\
&=
\frac{z^2}{2w}.
\end{align*}

Por lo tanto, la \textcolor{blue}{\textbf{función de oferta de bienes}} es:
\begin{equation}
  \boxed{Y^s(w) = 
  \frac{z^2}{2w} \label{eq:Y_s}
  }
\end{equation}
Con pendiente negativa y convexa respecto al salario real:
\[
Y^s_w
=
-\frac{z^2}{2w^2}<0, \quad Y^s_{ww} = \frac{z^2}{w^3}>0, \quad \forall w>0
\]
lo cual refleja que el incremento del salario real se traduce en un incremento en los costos laborales, lo cual reduce la demanda laboral y por lo tanto la producción de bienes.


La función de beneficios de la empresa es:
\begin{align*}
\pi(w)
&=Y^s(w)-wN^d(w)\\
&=\frac{z^2}{2w}
-w\left(\frac{z^2}{4w^2}\right)\\
&=\frac{z^2}{2w}-\frac{z^2}{4w}\\
&=
\boxed{
\frac{z^2}{4w}
}.
\end{align*}
Con pendiente negativa y convexa respecto al salario real:
\[
\pi_w
=
-\frac{z^2}{4w^2}<0, \quad \pi^s_{ww} = \frac{z^2}{2w^3}>0, \quad \forall w>0
\]

\subsubsection*{Oferta laboral en equilibrio general}

Finalmente, para obtener la función de oferta laboral, sustituimos
\[
\pi(w)=\frac{z^2}{4w}
\qquad\text{y}\qquad
G = T
\]
en la demanda de ocio del hogar:
\begin{align*}
\ell(w)
&=
\frac{\psi\left[w+\pi(w)-G\right]}
     {(1+\psi)w}
\end{align*}

Y así obtenemos la \textcolor{blue}{\textbf{función de demanda de ocio}}:
\begin{equation}
  \boxed{\ell(w) = 
  \frac{\psi}{(1+\psi)w}\left[w+\dfrac{z^2}{4w}-G \right]
},
\end{equation}

con la cual obtenemos la \textcolor{blue}{\textbf{función de oferta laboral}}:
\begin{equation}
  \boxed{N^s(w) = 
  1 - 
\frac{\psi}{(1+\psi)w}\left[w+\dfrac{z^2}{4w}-G \right]}.
\end{equation}

Desarrollando la expresión,
\[
\boxed{
N^s(w)
=
\frac{1}{1+\psi}
+\frac{\psi G}{(1+\psi)w}
-\frac{\psi z^2}{4(1+\psi)w^2}
}.
\]

La pendiente de esta función es
\begin{align*}
N^s_w
&=
-\frac{\psi G}{(1+\psi)w^2}
+\frac{\psi z^2}{2(1+\psi)w^3}\\[0.3cm]
&=
\frac{\psi(z^2-2Gw)}
     {2(1+\psi)w^3}
.
\end{align*}

En consecuencia, la oferta laboral no tiene pendiente positiva para
todo \(w>0\). Su pendiente es positiva solo cuando
\[
w<\frac{z^2}{2G}.
\]

Para saber si en equilibrio la oferta laboral tiene pendiente positiva, negativa o igual a cero, es necesario encontrar el salario real de equilibrio $w^*$.

Por otro lado, a partir de la segunda derivada observamos la curvatura en cada región:
\begin{align*}
N^s_{ww}
&=
\frac{2\psi G}{(1+\psi)w^3}
-\frac{3\psi z^2}{2(1+\psi)w^4}\\[0.3cm]
&=
\frac{\psi(4Gw - 3z^2)}
     {2(1+\psi)w^4}.
\end{align*}
de modo que para $w < \tfrac{3z^2}{4G}$ es cóncava y para $w > \tfrac{3z^2}{4G}$ es convexa. Como $\tfrac{1}{2} < \tfrac{3}{4}$, entonces la función de oferta laboral es cóncava en toda la región en que tiene pendiente positiva, y tiene un punto de inflexión en algún punto dentro de la región en que tiene pendiente negativa.

\subsubsection*{Demanda de bienes en equilibrio general}

La demanda de consumo privado se obtiene sustituyendo
\[
\pi(w)=\frac{z^2}{4w}
\qquad\text{y}\qquad
T=G
\]
en la demanda condicional de consumo:
\begin{align*}
C^d(w) =
\frac{w+\pi(w)-G}{1+\psi}
\end{align*}

Por lo tanto, la \textcolor{blue}{\textbf{función demanda de consumo}} es: 
\begin{equation}
  \boxed{C^d(w) = \frac{1}{1+\psi} \left[w+\dfrac{z^2}{4w}-G \right]}
\end{equation}

Su pendiente respecto del salario real es
\begin{align*}
C^d_w
=
\frac{1}{1+\psi}
\left(
1-\frac{z^2}{4w^2}
\right),
\end{align*}

Como
\[
N^d(w)=\frac{z^2}{4w^2},
\]
también puede escribirse como
\[
C^d_w
=
\frac{1-N^d(w)}{1+\psi}.
\]
Por tanto, en un equilibrio en el que \(N^{d*}<1\), la demanda de consumo privado tiene pendiente positiva respecto del salario real.

Por otro lado, su curvatura es:
$$
C^d_{ww} = \frac{1}{1+\psi}
\left(
\frac{z^2}{2w^3}
\right) > 0, \quad  \forall w>0,
$$
de modo que es convexa respecto al salario real.

Por último, la \textcolor{blue}{\textbf{función de demanda agregada}} es:
\[
Y^d(w)=C^d(w)+G.
\]
Sustituyendo la función de consumo,
\begin{equation}
\boxed{Y^d(w)
=
\frac{1}{1+\psi}
\left[w+\dfrac{z^2}{4w}+\psi G \right]
}.
\end{equation}

De esta forma, la economía se representa por las siguientes funciones:
\begin{equation}
  \boxed{
    \begin{aligned}
      N^d(w) &= \frac{z^2}{4w^2}
        & \text{(demanda laboral)} \\[2mm]
      N^s(w) &= \frac{1}{1+\psi}+\frac{\psi G}{(1+\psi)w}-\frac{\psi z^2}{4(1+\psi)w^2}
        & \text{(oferta laboral)} \\[2mm]
      Y^s(w) &= \frac{z^2}{2w}
        & \text{(oferta agregada de bienes)} \\[2mm]
      Y^d(w) &= \frac{1}{1+\psi}\left[w+\frac{z^2}{4w}+\psi G\right]
        & \text{(demanda agregada de bienes)}
    \end{aligned}
  }
\end{equation}

Alternativamente, podemos representar a la economía por medio de las funciones que determinan cada una de las variables macroeconómicas:
\begin{equation}
  \boxed{
    \begin{aligned}
      N^d(w) &= \frac{z^2}{4w^2}
        & \text{(demanda laboral)} \\[2mm]
      N^s(w) &= \frac{1}{1+\psi}+\frac{\psi G}{(1+\psi)w}-\frac{\psi z^2}{4(1+\psi)w^2}
        & \text{(oferta laboral)} \\[2mm]
      Y^s(w) &= \frac{z^2}{2w}
        & \text{(oferta agregada de bienes)} \\[2mm]
      Y^d(w) &= \frac{1}{1+\psi}\left[w+\frac{z^2}{4w}+\psi G\right]
        & \text{(demanda agregada de bienes)} \\[2mm]
      \pi(w) &= \frac{z^2}{4w}
        & \text{(beneficios empresariales)} \\[2mm]
      \ell(w) &= \frac{\psi}{(1+\psi)w}\left[w+\frac{z^2}{4w}-G\right]
        & \text{(demanda de ocio)} \\[2mm]
      C^d(w) &= \frac{1}{1+\psi}\left[w+\frac{z^2}{4w}-G\right]
        & \text{(demanda de consumo privado)}
    \end{aligned}
  }
\end{equation}

El salario real de equilibrio \(w^*\) es el precio relativo que cumple con el vaciado de mercados. Es decir, es el nivel salarial tal que:
\begin{align*}
  N^s(w^*)&=N^d(w^*), \\
  Y^s(w^*)&=Y^d(w^*).
\end{align*}

Como solo tenemos dos mercados, por ley de Walras sabemos que basta con encontrar el salario de equilibrio que vacía uno de los dos mercados para garantizar que el otro mercado se vacíe automáticamente. 

El siguiente inciso nos pide \textbf{encontrar ese precio de equilibrio $w^*$ y las cantidades de equilibrio asociadas}. Es decir, encontrar el \textbf{equilibrio competitivo}.



\subsection*{(c) Solución explícita del equilibrio competitivo}
En el inciso anterior encontramos las elecciones óptimas del hogar y de la empresa como funciones del salario real y de las variables exógenas. Con estas funciones, ahora es posible encontrar la asignación del equilibrio competitivo: tan solo necesitamos encontrar el salario real $w^*$ de equilibrio para sustituirlo en las funciones de cada variable y así encontrar las cantidades de equilibrio.

Para encontrar el salario real de equilibrio, utilizamos la condición de vaciado del mercado de trabajo:
$$
N^s(w^*) = N^d(w^*)
$$

Sustituyendo las funciones de oferta y demanda laboral,
\[
  \frac{1}{1+\psi}+\frac{\psi G}{(1+\psi)w}-\frac{\psi z^2}{4(1+\psi)w^2}
  =\frac{z^2}{4w^2}.
\]
Multiplicando ambos lados por $(1+\psi)w$,
\[
  w + \psi G -\frac{\psi z^2}{4w}
  =\frac{(1+\psi)z^2}{4w}.
\]
Multiplicando ahora por $4w$,
\[
  4w^2 + 4\psi Gw - \psi z^2 = (1+\psi)z^2.
\]
Reuniendo todos los términos,
\[
  4w^2+4\psi Gw-(1+2\psi)z^2=0.
\]
Utilizando la fórmula general, vemos que esta ecuación produce dos raíces
\[
  w
  =\frac{-4\psi G\pm
  \sqrt{16\psi^2G^2+16(1+2\psi)z^2}}{8}.
\]
Como el salario debe ser positivo, se descarta la raíz negativa. Por lo tanto,
\[
  \boxed{
  w^*
  =\frac{-\psi G+\sqrt{\psi^2G^2+(1+2\psi)z^2}}{2}
  }.
\]
Sustituyendo en la función de demanda laboral, obtenemos el empleo de equilibrio:
\[
  \boxed{
  N^*
  =\left[
  \frac{\psi G+\sqrt{\psi^2G^2+(1+2\psi)z^2}}
  {(1+2\psi)z}
  \right]^2
  }.
\]
El ocio de equilibrio es
\[
  \boxed{
  \ell^*
  =1-\left[
  \frac{\psi G+\sqrt{\psi^2G^2+(1+2\psi)z^2}}
  {(1+2\psi)z}
  \right]^2
  }.
\]
La producción de equilibrio es
\[
  Y^*=z\sqrt{N^*}
  =\boxed{
  \frac{\psi G+\sqrt{\psi^2G^2+(1+2\psi)z^2}}
  {1+2\psi}
  }.
\]
El consumo se obtiene del vaciado del mercado de bienes:
\[
  C^*=Y^*-G,
\]
por lo que
\[
  \boxed{
  C^*
  =\frac{\sqrt{\psi^2G^2+(1+2\psi)z^2}-(1+\psi)G}
  {1+2\psi}
  }.
\]
Finalmente,
\[
  \boxed{G = T^*}.
\]

Así econtramos el equilibrio competitivo.\footnote{Noten que resolver el equilirbio competitio es encontrar la asignación y el salario real de equilibrio en términos únicamente del parámetro $\psi$ y de las variables exógenas $G$ y $z$.}

Por lo tanto, el \textbf{equilibrio competitivo} de la economía es:
\begin{equation}
  \boxed{
    \begin{aligned}
      w^* &= \frac{-\psi G+\sqrt{\psi^2G^2+(1+2\psi)z^2}}{2}
        & \text{(salario real de equilibrio)} \\[3mm]
      N^* &= \left[\frac{\psi G+\sqrt{\psi^2G^2+(1+2\psi)z^2}}{(1+2\psi)z}\right]^2
        & \text{(empleo de equilibrio)} \\[3mm]
      \ell^* &= 1-\left[\frac{\psi G+\sqrt{\psi^2G^2+(1+2\psi)z^2}}{(1+2\psi)z}\right]^2
        & \text{(ocio de equilibrio)} \\[3mm]
      Y^* &= \frac{\psi G+\sqrt{\psi^2G^2+(1+2\psi)z^2}}{1+2\psi}
        & \text{(producción de equilibrio)} \\[3mm]
      C^* &= \frac{\sqrt{\psi^2G^2+(1+2\psi)z^2}-(1+\psi)G}{1+2\psi}
        & \text{(consumo de equilibrio)} \\[3mm]
      \pi^* &= \frac{\psi G+\sqrt{\psi^2G^2+(1+2\psi)z^2}}{2(1+2\psi)}
        & \text{(beneficios de equilibrio)} \\[3mm]
      T^* &= G
        & \text{(impuesto de equilibrio)}
    \end{aligned}
  }
\end{equation}
En equilibrio, el precio $w^*$ es tal que: la oferta y demanda laboral coinciden, generando un nivel de empleo de equilibrio $N^*$; la oferta y demanda agregada de bienes y servicios coinciden, generando un nivel de producción de equilibrio $Y^*$.


\subsubsection*{Análisis del equilibrio competitivo}
Como vimos en el inciso anterior, la oferta de trabajo no tiene pendiente positiva para todo \(w>0\), de modo que para un rango salarial determinado, el efecto ingreso puede dominar al efecto sustitución en la elección de consumo-ocio del hogar. Analizaremos si en el equlibrio y en regiones cercanas a este, la oferta laboral muestra un efecto ingreso pronunciado.

La pendiente es positiva solo cuando
\[
w<\frac{z^2}{2G}.
\]
En equilibrio, el salario real es igual a:
$$
w^* = \frac{-\psi G+\sqrt{\psi^2G^2+(1+2\psi)z^2}}{2}
$$

Sustituyendo,
\[
w^*<\frac{z^2}{2G}
\quad\Longleftrightarrow\quad
-\psi G+\sqrt{\psi^2G^2+(1+2\psi)z^2}<\frac{z^2}{G}
\quad\Longleftrightarrow\quad
\sqrt{\psi^2G^2+(1+2\psi)z^2}<\frac{z^2+\psi G^2}{G}.
\]
Como ambos lados son positivos, podemos elevar al cuadrado sin alterar el sentido de la desigualdad. Multiplicando también por $G^2$,
\[
G^2\left[\psi^2G^2+(1+2\psi)z^2\right]<\left(z^2+\psi G^2\right)^2.
\]
Desarrollando ambos lados,
\[
\psi^2G^4+(1+2\psi)z^2G^2 <
z^4+2\psi G^2z^2+\psi^2G^4.
\]
Cancelando el término $\psi^2G^4$ de ambos lados y reagrupando,
\[
(1+2\psi)z^2G^2-2\psi G^2z^2<z^4
\quad\Longleftrightarrow\quad
z^2G^2<z^4
\quad\Longleftrightarrow\quad
G^2<z^2.
\]
Es decir,
\[
w^*<\frac{z^2}{2G}
\quad\Longleftrightarrow\quad
G<z
\]
(usando que $G,z>0$). Como el ejercicio supone $0<G<z$, esta condición se satisface, por lo que
\[
\boxed{w^*<\frac{z^2}{2G}}.
\]

En consecuencia, bajo el supuesto $0<G<z$, la oferta de trabajo tiene \textbf{pendiente positiva en el equilibrio}: el efecto sustitución domina al efecto ingreso en las cercanías de $w^*$. Intuitivamente, un nivel de gasto público suficientemente bajo hace que la necesidad que tiene el gobierno de aumentar los impuestos para financiar su gasto sea lo suficientemente baja para no generar un efecto ingreso negativo grande sobre el ingreso disponible del hogar. Si el gasto público creciera de $G$ a $G'$ suficientemente alto, de tal forma que
\[
w^*>\frac{z^2}{2G'},
\]
entonces los impuestos que el gobierno se vería en la necesidad de cobrar serían tan altos que el efecto ingreso negativo le ganaría al efecto sustitución, generando una pendiente negativa en la oferta laboral: una mejora en el salario le permitiría al hogar financiar un poco más de ocio, dado que el salario es muy bajo respecto a las necesidades que tiene para cubrir la caída de su ingreso disponible.


Otra propiedad del equilibrio competitivo que encontramos es que la condición $G<z$ garantiza una solución interior para el problema del hogar. 
Para ver esto, veamos que $C^*>0$ si y solo si
\[
  \sqrt{\psi^2G^2+(1+2\psi)z^2}>(1+\psi)G
\]
Y dado que
\[
  \psi^2G^2+(1+2\psi)z^2-(1+\psi)^2G^2
  =(1+2\psi)(z^2-G^2)>0,
\]
entonces $C^*>0$. Asimismo,
\[
  \left[(1+2\psi)z-\psi G\right]^2
  -\left[\psi^2G^2+(1+2\psi)z^2\right]
  =2\psi(1+2\psi)z(z-G)>0.
\]
Esto implica
\[
  \psi G+\sqrt{\psi^2G^2+(1+2\psi)z^2}<(1+2\psi)z,
\]
de modo que $0<N^*<1$ y $0<\ell^*<1$, lo que descarta soluciones de equina.

\subsection*{(d) Vaciado del mercado de bienes}

En el inciso anterior, utilizamos solo el vaciado del mercado de trabajo para encontrar el salario real de equilibrio $w^*$, de modo que a dicho salario el mercado laboral está en equilibrio por construcción. En este inciso, demostraremos que basta con que este mercado esté en equilibrio para que el mercado de bienes también lo esté.


En equilibrio, la restricción presupuestaria del hogar se satisface con igualdad:
\[
  C^*=w^*N^{s*}+\pi^*-T^*.
\]
La identidad de beneficios de la empresa es
\[
  \pi^*=Y^*-w^*N^{d*}.
\]
El presupuesto del gobierno satisface $G = T^*$ y el mercado de trabajo se vacía, por lo que
\[
  N^{s*}=N^{d*}=N^*.
\]
Sustituyendo estas tres condiciones en el presupuesto del hogar,
\[
  C^*
  =w^*N^*+Y^*-w^*N^*-G,
\]
\[
  C^*=Y^*-G.
\]
Por lo tanto,
\[
  \boxed{Y^*=C^*+G}.
\]
Así, si el mercado de trabajo se vacía, el mercado de bienes se vacía automáticamente (ley de Walras).

\subsection*{(e) Problema del planificador social}

El planificador social busca maximizar el bienestar del consumidor dadas las restricciones materiales de la economía:
\begin{align*}
  \max_{C,\ell}
  &\left\{\log C+\psi\log \ell \right\} \\
  \text{s.a.} \quad &C+G\leq z\sqrt N, \\
  &N + \ell \leq 1
\end{align*}
Sustituyendo la segunda restricción en la función de utilidad, reducimos las dos restricciones a una.

El Lagrangeano asociado entonces es:
\[
  \mathcal L(C, N, \mu)
  =\log C+\psi\log(1-N)
  +\mu\left(z\sqrt N-C-G\right).
\]
Las condiciones de primer orden son
\[
  \frac{1}{C}-\mu=0,
\]
\[
  -\frac{\psi}{1-N}+\mu\frac{z}{2\sqrt N}=0,
\]
\[
  z\sqrt N-C-G=0.
\]
De la primera condición, $\mu=1/C$. Al introducirla en la segunda,
\[
  \frac{\psi}{1-N}=\frac{z}{2C\sqrt N}.
\]
Por lo tanto,
\[
  C=\frac{z(1-N)}{2\psi\sqrt N}.
\]
Sustituyendo en la tercera condición, se obtiene que
\[
  z\sqrt N-G=\frac{z(1-N)}{2\psi\sqrt N}.
\]
Multiplicando por $2\psi\sqrt N$,
\[
  2\psi zN-2\psi G\sqrt N=z-zN.
\]
Reuniendo términos,
\[
  (1+2\psi)zN-2\psi G\sqrt N-z=0.
\]
Esta ecuación es cuadrática en $\sqrt N$. Su raíz positiva es
\[
  \sqrt{N^*}
  =\frac{2\psi G+\sqrt{4\psi^2G^2+4(1+2\psi)z^2}}
  {2(1+2\psi)z},
\]
es decir,
\[
  \sqrt{N^*}
  =\frac{\psi G+\sqrt{\psi^2G^2+(1+2\psi)z^2}}
  {(1+2\psi)z}.
\]
Por consiguiente, el planificador elige exactamente
\[
  N^*
  =\left[
  \frac{\psi G+\sqrt{\psi^2G^2+(1+2\psi)z^2}}
  {(1+2\psi)z}
  \right]^2,
\]
\[
  C^*
  =\frac{\sqrt{\psi^2G^2+(1+2\psi)z^2}-(1+\psi)G}
  {1+2\psi},
\]
y $\ell^*=1-N^*$. Estas son las mismas cantidades obtenidas en el equilibrio competitivo.

La razón es que el planificador social asigna recursos de tal forma que la relación marginal de sustitución del hogar sea igual a la productividad marginal del trabajo. Esto, como lo sabemos, también caracteriza a la asignación del equilibrio competitivo.

\subsection*{(f) Construcción de los tres diagramas}

\subsubsection*{Mercado laboral}
El mercado laboral se grafica en el espacio $(N,w)$. Las funciones a graficar son las curvas de oferta y demanda laboral:
\begin{equation}
  \boxed{
    \begin{aligned}
      N^d(w) &= \frac{z^2}{4w^2}
        & \text{(demanda laboral)} \\[2mm]
      N^s(w) &= \frac{1}{1+\psi}+\frac{\psi G}{(1+\psi)w}-\frac{\psi z^2}{4(1+\psi)w^2}
        & \text{(oferta laboral)}
    \end{aligned}
  }
\end{equation}
El gráfica debe incluir notaciones para indicar el equilibrio inicial.

\subsubsection*{Mercado de bienes}
El mercado de bienes y servicios se grafica en el espacio $(Y,w)$. Las funciones a graficar son las curvas de oferta y demanda agregada:
\begin{equation}
  \boxed{
    \begin{aligned}
      Y^s(w) &= \frac{z^2}{2w}
        & \text{(oferta agregada de bienes)} \\[2mm]
      Y^d(w) &= \frac{1}{1+\psi}\left[w+\frac{z^2}{4w}+\psi G\right]
        & \text{(demanda agregada de bienes)}
    \end{aligned}
  }
\end{equation}

\subsubsection*{Asignación en el plano $(\ell,C)$}
Esta es la gráfica que une la FPP con las curvas de indiferencia del agente y su restricción presupuestaria. Representa de forma gráfica la asignación del equilibrio competitivo y cambios en este equilibrio a partir un choque exógeno.



\subsection*{(g) Choque de gasto público}

Este tipo de incisos buscan analizar el impacto de un choque exógeno sobre el equilibrio competitivo de la economía, así como sus canales de transmisión. Primero nos preguntamos: ¿qué sucede con las cantidades y precios de equilibrio en la economía?

Hay dos formas equivalentes de responder a esta pregunta: analíticamente y gráficamente. Analíticamente lo podemos hacer a través de derivar las variables en equilibrio respecto a la variación exógena. Gráficamente lo hacemos al analizar desplazamientos de y sobre las curvas en los mercados involucrados.

Comenzaremos con el estudio analítico.\footnote{Para efectos de examen, es recomendable que dominen la intuición gráfica, pues les tomará mucho menos tiempo que hacerlo de forma analítica y es un procedimiento equivalente. Para efectos de tareas y ejercicios de estudio, es recomendable hacer ambos.}

La respuesta del salario real es:
\[
  \frac{\partial w^*}{\partial G}
  =\frac{1}{2}
  \left[
  -\psi + 
  \frac{\psi^2G}{\sqrt{\psi^2G^2+(1+2\psi)z^2}}
  \right].
\]
Por lo tanto, 
\begin{align*}
  \frac{\partial w^*}{\partial G} <0 
  &\iff -\psi + 
  \frac{\psi^2G}{\sqrt{\psi^2G^2+(1+2\psi)z^2}} <0 \\
  &\iff 
  \frac{\psi^2G}{\sqrt{\psi^2G^2+(1+2\psi)z^2}} < \psi \\
  &\iff
  \psi G  < \sqrt{\psi^2G^2+(1+2\psi)z^2} \\
  &\iff
  \psi^2 G^2  < \psi^2G^2+(1+2\psi)z^2 \\
  &\iff
  0 < (1+2\psi)z^2 \\
\end{align*}
Como $z>0$ y $\psi > 0$, entonces $0 < (1+2\psi)z^2$ se cumple siempre. Por lo tanto, el salario real de equilibrio cae sin ambigüedad con aumentos en el gasto público:
\[
  \boxed{\frac{\partial w^*}{\partial G}<0}.
\]
La respuesta del empleo es
\[
  \frac{\partial N^*}{\partial G}
  =\frac{2\left[\psi G+\sqrt{\psi^2G^2+(1+2\psi)z^2}\right]}
  {(1+2\psi)^2z^2}
  \left[
  \psi+\frac{\psi^2G}{\sqrt{\psi^2G^2+(1+2\psi)z^2}}
  \right]>0.
\]
Por lo tanto, el empleo aumenta y el ocio disminuye:
\[
  \boxed{\frac{\partial N^*}{\partial G}>0},
  \qquad
  \boxed{\frac{\partial\ell^*}{\partial G}<0}.
\]

La respuesta de la producción es positiva
\[
\boxed{
  \frac{\partial Y^*}{\partial G}
  = \frac{
  1}{1+2\psi} \left[
  \psi+\frac{\psi^2G}{\sqrt{\psi^2G^2+(1+2\psi)z^2}}
  \right] >0
}.
\]


La respuesta del consumo es:
\[
  \frac{\partial C^*}{\partial G} 
  =
  \frac{
  1}{1+2\psi} \left[
  \dfrac{\psi^2G}{\sqrt{\psi^2G^2+(1+2\psi)z^2}}
  -(1+\psi)
  \right].
\]
Por lo tanto,
\begin{align*}
  \frac{\partial C^*}{\partial G} < 0 
  &\iff 
  \dfrac{\psi^2G}{\sqrt{\psi^2G^2+(1+2\psi)z^2}}
  -(1+\psi) < 0 \\
  &\iff 
  \psi^2G < (1+\psi)\sqrt{\psi^2G^2+(1+2\psi)z^2} \\
  &\iff 
  \psi^4 G^2 < (1+\psi)^2\left(\psi^2G^2+(1+2\psi)z^2\right) \\
  &\iff 
  \psi^4 G^2 - (1+\psi)^2\psi^2 G^2 < (1+\psi)^2(1+2\psi) z^2 \\
  &\iff 
  \psi^2 G^2\left[\psi^2 - (1+\psi)^2\right] < (1+\psi)^2(1+2\psi) z^2 \\
  &\iff 
  \psi^2 G^2\left[\psi^2 - (\psi^2 + 2 \psi + 1)\right] < (1+\psi)^2(1+2\psi) z^2 \\
  &\iff 
  -\psi^2 G^2 (1+2\psi) < (1+\psi)^2(1+2\psi) z^2 \\
  &\iff 
  0< \psi^2 G^2 (1+2\psi) + (1+\psi)^2(1+2\psi) z^2 \\
  &\iff 
  (1+2\psi)\left[\psi^2 G^2 + (1+\psi)^2 z^2\right] > 0.
\end{align*}
Como $z>0$ y $\psi > 0$, entonces $(1+2\psi)\left[\psi^2 G^2 + (1+\psi)^2 z^2\right] > 0$ se cumple siempre. Por lo tanto, el consumo cae sin ambigüedad con aumentos en el gasto público:
\[
  \boxed{\frac{\partial C^*}{\partial G} < 0.}
\]

Por último, como $Y^* = C^* + G$, entonces $C^* = Y^* - G$, de modo que si $\tfrac{\partial Y^*}{\partial G} >0$ y $\tfrac{\partial C^*}{\partial G} <0$, entonces
\[
\frac{\partial C^*}{\partial G} = \frac{\partial Y^*}{\partial G}  - 1 < 0 \implies  < \frac{\partial Y^*}{\partial G} < 1
\]
Es decir, el multiplicador fiscal es positivo, pero menor que uno.
\[
  \boxed{
  0<\frac{\partial Y^*}{\partial G}<1
  }.
\]

Después de este análisis gráfico, se analiza el choque y su propagación en la economía a través de las funciones que obtuvimos en incisos anteriores y sus gráficas asociadas. Es decir, la propagación a través del mercado laboral, del mercado de bienes y servicios, y del diagrama que muestra la asignación de equilibrio competitivo en el espacio $(\ell, C)$.

\paragraph{Interpretación.}

El mecanismo económico es el siguiente. El aumento en el gasto público requiere forzosamente ser financiado con un aumento de impuestos. Esto genera un choque negativo en el ingreso disponible de los hogares, lo cual se traduce en un efecto ingreso negativo que hace que el hogar se sienta más pobre con el nuevo nivel de consumo y se le generen incentivos para trabajar, de modo que aumenta su oferta laboral. Lo anterior genera un desplazamiento hacia la derecha de la curva de oferta en el mercado laboral:
\[
  \frac{\partial N^s(w)}{\partial G}
  =\frac{\psi}{(1 + \psi)w} > 0
\]
lo que genera presiones a la baja en el salario real. Formalmente, esto se mostró al ver que $\partial w^*/\partial G<0$. El salario real así se ajusta a la baja, lo que genera una aumento en la demanda de trabajo que se traduce como un movimiento sobre la curva de demanda laboral:
\[
  \frac{\partial N^d(w)}{\partial w}
  = -\frac{z^2}{2w^3} < 0
\]
Por lo tanto, el mercado laboral se ajusta a un nuevo equilibrio con menor nivel salarial pero mayor empleo, resultado que también se confirma analíticamente ($\partial N^*/\partial G>0$).

Este incremento en el empleo resulta en un incremento de la producción de bienes y servicios. Sin embargo, este incremento es \textit{menor} que el propio incremento del gasto público, pues $\partial Y^*/\partial G<1$. Además, el consumo privado de equilibrio \textit{cae} ante el choque ($\partial C^*/\partial G<0$), aunque en una magnitud menor a la del propio choque de gasto. Es decir, el ajuste endógeno del empleo y la producción amortigua, pero no revierte, la caída inicial en el ingreso disponible.

El incremento en el gasto público se traduce en un incremento en la demanda agregada de la economía, desplazando hacia la derecha la curva de demanda en el mercado de bienes y servicios:
$$
\frac{\partial Y^d(w)}{\partial G} = \frac{\psi}{1 + \psi} >0
$$
Por otro lado, el incremento en el nivel de empleo resulta en un incremento de la producción de bienes y servicios. Este nuevo equilibrio en el mercado laboral reduce el nivel de salario real, lo que genera un movimiento hacia la derecha sobre la curva de oferta agregada en el mercado de bienes y servicios:
$$
\frac{\partial Y^s(w)}{\partial w} = -\frac{z^2}{2w^2} <0
$$
Todo lo anterior resulta en un nivel de producción de equilibrio mayor, aunque con un gasto en consumo privado menor.\footnote{Toda esta interpretación descansa en el supuesto $0<G<z$, que como mostramos en el inciso anterior, asegura que el salario de equilibrio se ubique siempre en la región de pendiente positiva de la oferta laboral ($w^*<z^2/(2G)$). Cabe destacar que, dentro de este rango de $G$, un aumento del gasto público \textit{nunca} puede generar una caída del empleo de equilibrio: la derivada $\partial N^*/\partial G$ es estrictamente positiva para todo $G\in(0,z)$, sin importar cuán cerca esté $G$ de su cota superior $z$. Sin embargo, en el límite $G\to z^-$ se tiene $N^*\to1$ (y $\ell^*\to0$), de modo que el hogar destinaría la totalidad de su dotación de tiempo al trabajo. Esto revela que el supuesto $G<z$ no es solo una restricción técnica, sino la condición que garantiza que exista una asignación interior de ocio positivo; si el gasto público se elevara hasta igualar o superar a la productividad $z$, los hogares destinarían todo su tiempo a trabajar para poder recuperar en mayor medida posible el poder adquisitivo perdido por el nivel tan alto de gasto público.}

\subsection*{(h) Choque de productividad}

La respuesta del salario es
\[
  \boxed{
  \frac{\partial w^*}{\partial z}
  =\frac{(1+2\psi)z}
  {2\sqrt{\psi^2G^2+(1+2\psi)z^2}}>0
  }.
\]
La producción y el consumo responden de acuerdo con
\[
  \boxed{
  \frac{\partial Y^*}{\partial z}
  =\frac{z}{\sqrt{\psi^2G^2+(1+2\psi)z^2}}>0
  },
\]
\[
  \boxed{
  \frac{\partial C^*}{\partial z}
  =\frac{z}{\sqrt{\psi^2G^2+(1+2\psi)z^2}}>0
  }.
\]
La respuesta del empleo es:
\begin{align*}
  \frac{\partial N^*}{\partial z}
  = 2 \left[\frac{\psi G+\sqrt{\psi^2G^2+(1+2\psi)z^2}}{(1+2\psi)z}\right] \cdot \frac{\partial}{\partial z}\left[\frac{\psi G+\sqrt{\psi^2G^2+(1+2\psi)z^2}}{(1+2\psi)z}\right]
\end{align*}
Como $2 \left[\frac{\psi G+\sqrt{\psi^2G^2+(1+2\psi)z^2}}{(1+2\psi)z}\right] > 0$, entonces la dirección de respuesta del empleo depende del signo de $\frac{\partial}{\partial z}\left[\frac{\psi G+\sqrt{\psi^2G^2+(1+2\psi)z^2}}{(1+2\psi)z}\right]$.

Por regla del cociente,
\begin{align*}
  \frac{\partial}{\partial z}\left[\frac{\psi G+\sqrt{\psi^2G^2+(1+2\psi)z^2}}{(1+2\psi)z}\right] &= \frac{\dfrac{(2(1+2\psi)z)(1+2\psi)z}{2\sqrt{\psi^2G^2+(1+2\psi)z^2}} - (1+2\psi) \left( \psi G+\sqrt{\psi^2G^2+(1+2\psi)z^2} \right)}{(1+2\psi)^2 z^2} \\[4pt]
  &= \frac{\dfrac{(1+2\psi)^2 z^2}{\sqrt{\psi^2G^2+(1+2\psi)z^2}} - (1+2\psi) \left( \psi G+\sqrt{\psi^2G^2+(1+2\psi)z^2} \right)}{(1+2\psi)^2 z^2} \\[4pt]
  &= \frac{1}{\sqrt{\psi^2G^2+(1+2\psi)z^2}} - \frac{ \psi G+\sqrt{\psi^2G^2+(1+2\psi)z^2} }{(1+2\psi) z^2} \\[4pt]
  &= \frac{(1+2\psi)z^2 - \sqrt{\psi^2G^2+(1+2\psi)z^2}\left(\psi G+\sqrt{\psi^2G^2+(1+2\psi)z^2}\right)}{(1+2\psi)z^2\sqrt{\psi^2G^2+(1+2\psi)z^2}} \\[4pt]
  &= \frac{(1+2\psi)z^2 - \psi G\sqrt{\psi^2G^2+(1+2\psi)z^2} - \left[\psi^2G^2+(1+2\psi)z^2\right]}{(1+2\psi)z^2\sqrt{\psi^2G^2+(1+2\psi)z^2}} \\[4pt]
  &= \frac{-\psi^2G^2 - \psi G\sqrt{\psi^2G^2+(1+2\psi)z^2}}{(1+2\psi)z^2\sqrt{\psi^2G^2+(1+2\psi)z^2}} \\[4pt]
  &= \frac{-\psi G\left[\psi G+\sqrt{\psi^2G^2+(1+2\psi)z^2}\right]}{(1+2\psi)z^2\sqrt{\psi^2G^2+(1+2\psi)z^2}}.
\end{align*}
Sustituyendo de vuelta,
\begin{align*}
  \frac{\partial N^*}{\partial z}
  &= 2 \left[\frac{\psi G+\sqrt{\psi^2G^2+(1+2\psi)z^2}}{(1+2\psi)z}\right] \cdot \frac{-\psi G\left[\psi G+\sqrt{\psi^2G^2+(1+2\psi)z^2}\right]}{(1+2\psi)z^2\sqrt{\psi^2G^2+(1+2\psi)z^2}} \\[4pt]
  &= -\,\frac{2\psi G\left[\psi G+\sqrt{\psi^2G^2+(1+2\psi)z^2}\right]^2}{(1+2\psi)^2 z^3 \sqrt{\psi^2G^2+(1+2\psi)z^2}}.
\end{align*}


Por lo tanto, para todo $z>0$:
\[
\boxed{
  \frac{\partial N^*}{\partial z} < 0
}.
\]

Es decir, ante un aumento en la productividad $z$, el empleo de equilibrio \textbf{disminuye}. Esto se debe a que, \textbf{para este tipo de preferencias, el efecto ingreso domina al efecto sustitución}: al ser el trabajador más productivo, un mismo nivel de consumo se puede sostener trabajando menos horas, y en equilibrio el hogar decide convertir parte del mayor ingreso que genera esa productividad en ocio.

Nótese que si $G=0$, la derivada del empleo respecto de la productividad es cero: sin gasto público, el efecto ingreso desaparece del canal de $N^*$ y el empleo de equilibrio se vuelve insensible (localmente) a cambios en $z$.

\paragraph{Interpretación.} El incremento en la productividad resulta en un incremento en el producto marginal del trabajo: dado el nivel salarial inicial, el nivel de empleo produce una cantidad mayor que antes. Esto le genera incentivos a las empresas para aumentar su demanda laboral, dado que ahora el producto marginal del trabajo es mayor. Esto desplaza la curva de demanda laboral hacia la derecha:
\[
  \frac{\partial N^d(w)}{\partial z}
  =\frac{z}{2w^2}>0,
\]
lo que genera presiones al alza en el salario de equilibrio. Por otro lado, la mayor demanda laboral y el ajuste subsecuente al alza del salario genera un incremento en el ingreso laboral del hogar. Este \textbf{incremento en el ingreso laboral genera un efecto sustitución y un efecto ingreso condicionales a las preferencias del hogar}.
Para este caso, observemos que
\[
  \frac{\partial N^s(w)}{\partial z}
  =-\frac{\psi z}{2(1+\psi)w^2}<0,
\]
por lo que la curva de oferta laboral se desplaza hacia la izquierda. Esto indica que el efecto ingreso domina al efecto sustitución, de tal forma que, dado que el hogar más productivo ahora puede alcanzar mayores niveles de consumo con un mismo nivel de empleo, entonces reduce su oferta laboral.

En el mercado de bienes, ambas curvas se desplazan hacia arriba. Para la producción,
\[
  \frac{\partial Y^s(w)}{\partial z}=\frac{z}{w}>0,
\]
y para la demanda agregada,
\[
  \frac{\partial Y^d(w)}{\partial z}
  =\frac{z}{2(1+\psi)w}>0.
\]
El desplazamiento directo de la producción es mayor que el desplazamiento de la demanda agregada a un salario dado. El salario aumenta hasta restablecer la igualdad $Y=C+G$. Deben mostrarse tanto los desplazamientos de las dos curvas como el movimiento hacia un salario y una producción mayores.

En el plano $(\ell,C)$,
\[
  \frac{\partial[z\sqrt{1-\ell}-G]}{\partial z}
  =\sqrt{1-\ell}>0.
\]
La frontera de posibilidades se desplaza hacia arriba y se vuelve más inclinada en valor absoluto, pues
\[\frac{\partial}{\partial z}
  PMN 
  =
  \frac{\partial}{\partial z}
  \left[-\frac{z}{2\sqrt{1-\ell}}\right]
  =-\frac{1}{2\sqrt{1-\ell}}<0.
\]
La nueva tangencia se encuentra arriba y a la derecha de la original: aumentan el consumo y el ocio.

\subsection*{(i) Identificación empírica de los choques}

Las respuestas predichas por el modelo, manteniendo constantes las demás variables exógenas, son
\[
\begin{array}{c|cccccc}
  \hline
  \text{Choque} & C & Y & N & \ell & w & \pi \\
  \hline
  G\uparrow & - & + & + & - & - & + \\
  z\uparrow\;(G>0) & + & + & - & + & + & + \\
  \hline
\end{array}
\]

Un aumento simultáneo de producción, consumo y salario real es consistente, dentro de este modelo, con un choque positivo de productividad. Un aumento de producción acompañado por disminuciones del consumo y del salario real es consistente con un aumento del gasto público financiado mediante impuestos de suma fija. Los beneficios aumentan en ambos casos, por lo que su signo no permite distinguir los dos choques. 

Para este caso, el empleo sí ayuda a distinguirlos: aumenta con el choque fiscal y disminuye con el choque de productividad cuando $G>0$. Sin embargo, este depende de las preferencias del hogar, pues si estas fueran tales que el hogar decide aumentar su oferta laboral dado el incremento en la productividad, entonces el nivel de empleo de equilibrio aumentaría, al igual que con el choque de gasto público.

\end{document}
